\chapter{Lie Groups}

Throughout the chapter, K denotes either the valued field R of real numbers, or the valued field C of complex numbers, or a non-discrete complete ultrametric commutative field. We assume that K is of characteristic 0 from § 4 onwards, that K = R or C in § 6, that K is ultrametric in § 7. Unless otherwise mentioned, all the manifolds, all the algebras and all the vector spaces considered are over K. Recall that, when we speak of a manifold of class \(C^{r}\), r ∈ N\_K, that is r = ω if K ≠ R and l ≤ r ≤ ω if K = R.

The conventions on norms, normable spaces and normed spaces are those of Differentiable and Analytic Manifolds, R.

Recall that a normable algebra over K is a (not necessarily associative) algebra A over K, with a topology T which possesses the following properties:

\begin{enumerate}
\def\labelenumi{(\arabic{enumi})}
\item
  \(\mathcal{T}\) can be defined by a norm;
\item
  the mapping \((x,y)\mapsto xy\) of \(\mathbf{A}\times \mathbf{A}\) into \(\mathbf{A}\) is continuous.
\end{enumerate}

The group of bicontinuous automorphisms of A is denoted by \(\mathrm{Aut(A)}\) . Every finite-dimensional algebra over \(K\) is a normable algebra with the canonical topology. A normed algebra over \(K\) is an algebra A over \(K\) with a norm such that \(\| xy\| \leqslant \| x\|\| y\|\) for all \(x,y\) in A; the algebra A with the topology defined by this norm is a normable algebra. If A is a normable algebra, there exists a norm on A defining its topology and making A into a normed algebra.

If G is a group, \(e_{G}\) , or simply e, denotes the identity element of G. For \(g \in G\) , \(\gamma(g)\) , \(\tilde{\delta}(g)\) and \(\operatorname{Int}(g)\) denote the mappings \(g' \mapsto gg'\) , \(g' \mapsto g'g^{-1}\) and \(g' \mapsto gg'g^{-1}\) of G into G. If f is a mapping of G into a set E, f denotes the mapping \(g \mapsto f(g^{-1})\) of G into E.

